\documentclass[10pt]{article}
\usepackage{setspace}
\usepackage{geometry}
\geometry{margin=0.7in}
\usepackage{graphicx}
\usepackage[hidelinks]{hyperref}
\setstretch{1.05}
\pagestyle{empty}

\begin{document}

\begin{flushright}
\textbf{Zhongli Wang} \\
Phone: +49 160 9541 6015 \\
Email: \href{mailto:nkwangzhongli@gmail.com}{nkwangzhongli@gmail.com} \\
Website: \href{https://sites.google.com/view/zhongli-wang/}{sites.google.com/view/zhongli-wang} \\
Astastrasse 22 \\
33617 Bielefeld \\
Germany \\
\end{flushright}

\noindent
\textbf{Dr Fran\c{c}ois Lafond} \\
Smith School of Enterprise and the Environment \\
School of Geography and the Environment \\
University of Oxford \\
Manor Road Building, Manor Road \\
Oxford OX1 3UQ, United Kingdom \\

\noindent
\textbf{Application for Postdoctoral Research Associate in the Economics of Technological Change (Vacancy 189392)}\\

\vspace{0.25cm}

\noindent Dear Dr Lafond,\\[0.15cm]

I am applying for the Postdoctoral Research Associate position in the Economics of Technological Change. I completed my PhD in Economics at Bielefeld University and Universit\`a Cattolica del Sacro Cuore in September 2026 through the Marie Sk{\l}odowska-Curie EPOC programme. My research combines complexity economics, industrial organisation, quantitative modelling, and artificial intelligence, closely matching the Oxford Martin Programme's objective of explaining and forecasting how technologies improve and diffuse.

My dissertation, \emph{Strategic Behavior and Economic Dynamics in ICT Markets: An Agent-Based Framework}, examines technological investment and market evolution under network effects, congestion, heterogeneous users, and regulation. In my job-market paper, firms adjust service-capacity investment while adoption responds to quality, prices, congestion, and peer influence. I build these models and reproducible Python pipelines, and my current work compares simulated method of moments with neural-network surrogates to estimate an ICT-industry model. These projects provide a strong basis for studying feedback among technology characteristics, costs, investment, and diffusion, and for testing whether structural information improves forecasts beyond historical trends.

I would contribute in three connected areas: developing a transparent technology database through reproducible collection, harmonisation, and validation; applying complex-systems models, econometrics, simulation, and machine learning to identify drivers of progress and assess predictions; and using LLMs for scalable data discovery and extraction. In \emph{Agentic Workflows for Economic Research}, my coauthors and I designed modular LLM-agent workflows using Microsoft AutoGen, tool use, multimodal inputs, and human oversight. I also work with LangChain, the OpenAI API, and Claude Code. This experience has made provenance, traceability, and quality control central to my research practice.

I have contributed to research from question formulation through publication. \emph{Polytope Fraud Theory}, published in the \emph{International Review of Financial Analysis} in 2025, translated qualitative warning signs into measurable indicators and evaluated them using supervised machine learning. My projects have also required cross-institutional coordination, computational documentation, international presentations, and workshop organisation. I would bring the same structured approach to publication leadership, research-assistant coordination, and partner collaboration.

Working across INET's Complexity Economics group, the Oxford Martin School, and the Smith School would connect my research on investment, adoption, and AI with a comparative programme spanning clean technologies and artificial intelligence. I would welcome the opportunity to contribute rigorous data, testable methods, and peer-reviewed research.

Thank you for considering my application. I look forward to the possibility of discussing my contribution to the team.

\vspace{0.2cm}

\noindent Sincerely,\\[0.3cm]
\noindent\includegraphics[width=2.3cm,keepaspectratio]{../Wang-sig.jpg}\\[-0.2cm]
\noindent Zhongli Wang

\end{document}
